\begin{helper}
Fix a finite terminal coordinate set \(U\), a number \(p\) with
\(0\le p\le 1/2\), and a fixed blue set \(B\subseteq U\) with
\(|B|=u_B\).  Let \(\mathcal B\) be a finite branch-label set with traces
\(\operatorname{blueTrace}(\beta)\), let \(\mathcal T\) be a finite set of
transcripts \(\tau=(P_\tau,A_\tau,Z_\tau)\), let \(J\) be a fixed red label,
and let \(R_{\beta,J}\) be the red support attached to \((\beta,J)\).  Suppose
there is a realized-cell predicate \(\operatorname{Realized}(\beta,\tau)\),
an assignment-compatibility predicate \(\operatorname{Comp}(A,\beta,\tau)\),
and a selected-present-sibling predicate
\(\operatorname{Sib}(A,\beta,\tau,e)\).

Assume the following reduced hypotheses, and no ambient history or terminal
product-law data.  For every realized \((\beta,\tau)\) with
\(\tau\in\mathcal T\), the realized-cell geometry gives
\[
|R_{\beta,J}|=u_R,\quad R_{\beta,J}\subseteq U,\quad
R_{\beta,J}\text{ red},\quad
P_\tau,A_\tau\subseteq U,\quad P_\tau\cap A_\tau=\emptyset,
\]
\[
B\cup R_{\beta,J}\subseteq P_\tau,\qquad Z_\tau\subseteq A_\tau,
\]
and for every blue \(e\in U\),
\[
e\in\operatorname{blueTrace}(\beta)\Longleftrightarrow e\in P_\tau,
\qquad
e\notin\operatorname{blueTrace}(\beta)\Longleftrightarrow e\in A_\tau .
\]
Also assume the span equivalence
\[
\operatorname{Comp}(A,\beta,\tau)\ \text{or}\
\exists e\in Z_\tau,\operatorname{Sib}(A,\beta,\tau,e)
\]
is equivalent, for every \(A\subseteq U\) and realized \((\beta,\tau)\), to
\[
P_\tau\subseteq A,\qquad
A\cap(A_\tau\setminus Z_\tau)=\emptyset,\qquad
\forall\text{ blue }e\in U\setminus Z_\tau,
  e\in\operatorname{blueTrace}(\beta)\Longleftrightarrow e\in A .
\]
Finally assume span disjointness: no \(A\subseteq U\) satisfies the left-hand
span condition for two distinct realized pairs.

Let
\[
L=\{(\beta,\tau):\tau\in\mathcal T,
      \operatorname{Realized}(\beta,\tau)\}.
\]
For \(\ell=(\beta,\tau)\in L\), let \(C_\ell\) be the set of
\(A\subseteq U\) satisfying the three displayed released-cylinder conditions
above.  Then
\[
\sum_{\ell=(\beta,\tau)\in L}
  p^{|P_\tau|}(1-p)^{|A_\tau|-|Z_\tau|}
\le
\sum_{A\in \bigcup_{\ell\in L} C_\ell}
  p^{|A|}(1-p)^{|U|-|A|}.
\]
\end{helper}
\begin{proof}
Fix a realized leaf \(\ell=(\beta,\tau)\in L\), and write
\(\tau=(P_\tau,A_\tau,Z_\tau)\).  Set
\[
D_\tau=A_\tau\setminus Z_\tau .
\]
The realized-cell geometry gives \(P_\tau,A_\tau\subseteq U\) and
\(P_\tau\cap A_\tau=\emptyset\).  Since \(D_\tau\subseteq A_\tau\), the
present set \(P_\tau\) and the absent set \(D_\tau\) are disjoint.  Because
\(Z_\tau\subseteq A_\tau\), deleting \(Z_\tau\) from \(A_\tau\) gives
\[
|D_\tau|=|A_\tau|-|Z_\tau|.
\]

We next check that the blue-trace part of \(C_\ell\) imposes no additional
product factors beyond \(P_\tau\) present and \(D_\tau\) absent.  Let \(e\in U\)
be blue.  The geometry gives
\[
e\in\operatorname{blueTrace}(\beta)\Longleftrightarrow e\in P_\tau
\]
and
\[
e\notin\operatorname{blueTrace}(\beta)\Longleftrightarrow e\in A_\tau .
\]
If \(e\notin Z_\tau\), the released-cylinder trace condition applies.  If
\(e\in P_\tau\), then \(e\in\operatorname{blueTrace}(\beta)\), so the trace
condition forces \(e\in A\).  If \(e\in A_\tau\setminus Z_\tau\), then
\(e\notin\operatorname{blueTrace}(\beta)\), so the trace condition forces
\(e\notin A\).  If \(e\in Z_\tau\), the coordinate is deliberately released:
its absent \((1-p)\)-factor has been removed from the residual cylinder, and
the trace condition imposes no value on it.  Thus selected blue absences are
free residual coordinates, while unselected absent blue coordinates remain in
the residual mass.  Conversely, for every unselected blue \(e\), membership in
\(P_\tau\) and \(A_\tau\) is exactly membership and nonmembership in the trace
by the two displayed equivalences, so the trace condition adds no third class
of fixed blue coordinates.

Thus \(C_\ell\) is precisely the complete terminal product cylinder with
coordinates in \(P_\tau\) fixed present and coordinates in \(D_\tau\) fixed
absent.  The free coordinates are the elements of
\(U\setminus(P_\tau\cup D_\tau)\).  Summing over each free coordinate
independently gives a factor \(p+(1-p)=1\); therefore
\[
\sum_{A\in C_\ell}p^{|A|}(1-p)^{|U|-|A|}
=p^{|P_\tau|}(1-p)^{|D_\tau|}
=p^{|P_\tau|}(1-p)^{|A_\tau|-|Z_\tau|}.
\]

It remains to pass from one leaf to all realized leaves.  Suppose
\(A\subseteq U\) lies in \(C_\ell\cap C_{\ell'}\) for two realized leaves
\(\ell=(\beta,\tau)\) and \(\ell'=(\beta',\tau')\).  By the span equivalence,
\(A\) satisfies the compatible-or-selected-sibling span condition for both
\((\beta,\tau)\) and \((\beta',\tau')\).  If \(\ell\ne\ell'\), this contradicts
the span-disjointness hypothesis.  Hence the cylinders \(C_\ell\) are
pairwise disjoint.  Summing the single-leaf identity over \(\ell\in L\), and
using pairwise disjointness to identify the sum of the cylinder masses with
the mass of their union, yields the asserted inequality.
\end{proof}
